\documentclass[a4paper, 12pt]{article}

%%%%%%%%%%%%%%%%%%%%%%%%%%%%%%%%%%%%%%%%%%%%%%%%%%%%%%%%%%%%%%%%%%%%%%%%%%%%%%%%

% Math, plots and sketches
\usepackage{amsmath, mathtools, amsfonts}
\usepackage{tkz-euclide}
\usepackage{pgfplots}
\pgfplotsset{compat=1.18}

% Images
\usepackage{graphicx}

% Allow line-breaking URLs/paths in text
\usepackage{url}
\def\UrlBreaks{\do\/\do-\do\_\do\#\do\0\do\1\do\2\do\3\do\4\do\5\do\6\do\7\do\8\do\9}

% Formatting
\usepackage[left=1in, right=1in, top=1.2in, bottom=1.2in]{geometry}
\usepackage{dirtytalk} % Proper quotation marks
\usepackage{float} % Make it possible for tables (and images) to be placed in an explicit location (using [H])


% Colors for highlighting and coloring tables
\usepackage{color, soul}
\usepackage{xcolor}

% Multicolumn formatting
% In case that's the format we want
%\usepackage{multicols}

% Custom labeled lists
\usepackage{enumerate}

% Clickable references
\usepackage{hyperref}
% Source - https://tex.stackexchange.com/a/51349
% Posted by Canageek, modified by community. See post 'Timeline' for change history
% Retrieved 2026-03-23, License - CC BY-SA 4.0
\hypersetup{
  colorlinks   = true, %Colours links instead of ugly boxes
  urlcolor     = blue, %Colour for external hyperlinks
  linkcolor    = blue, %Colour of internal links
  citecolor   = red %Colour of citations
}

% Comments
\usepackage{comment}
\usepackage{todonotes}

% References
\usepackage{biblatex}
\addbibresource{references.bib}

%%%%%%%%%%%%%%%%%%%%%%%%%%%%%%%%%%%%%%%%%%%%%%%%%%%%%%%%%%%%%%%%%%%%%%%%%%%%%%%%

\title{Counting Crows - Group 2b}

\author{ Kian Hamidi, Alex Janse, Jaël Kwakkel, Antonio Sato }

\date{ \today }

%%%%%%%%%%%%%%%%%%%%%%%%%%%%%%%%%%%%%%%%%%%%%%%%%%%%%%%%%%%%%%%%%%%%%%%%%%%%%%%%

\begin{document}
\maketitle

\section{Project Description}\label{sec:proj_desc}
% Short description of our project and its features
The behaviour of flocks of birds, schools of fish and other coordinated animal
group motion can be modelled using a simple set of steering behviours. This
project aims to implement and extend the behaviour of a flock of birds, in
specific a murder of crows, to run for 100,000 birds at a real-time simulation
rate of 60+ Hz. This consists of a simulation and a visualisation, both of
which are executed locally in a web application compiled to Web Assembly. This
report describes what software and frameworks have been used for the creation
of the simulation, how the simulation was implemented and what choices have
been made to create a real-time simulation for a large amount of birds. Finally
the problems we encountered while implementing this simulation are described,
before giving the results we have achieved. \todo{Update this section if we
    decide to add or remove chapters.}

\section{Software used}\label{sec:software}
% Provide details of software used (Threejs and WebGPU)
This chapter describes the choice of programming language, framework and
utilization of the GPU for both the simulation and visualization of the model.

\subsection{Simulation}\label{subsec:simulation}
% WebGPU, js, etc
The simulation is performed using a combination of Javascript on the CPU side,
and WGSL (WebGPU Shading Language) on the GPU side. The JavaScript code is
responsible for initializing GPU resources, dispatching compute workloads, and
synchronizing simulation parameters. WGSL is used to define a compute shader
that executes directly on the GPU, allowing parallel calculation of boid
behavior.

The advantages of using WGSL lied mostly in the performance gain. By allowing
parallel computation on the GPU, we were able to scale the simulation to
100,000+ agents, which is likely more than we could have achieved on CPU only.
Another advantage was the fine-grained control. By allowing us to explicitely
define buffers and memory layout, we were able to implement optimizations such
as tight memory packing, and reading back to the CPU in a structured way.

The main limitation we ran into when working with WGSL was the difficulty of debugging the program. Computation
happens in parallel, and variable inspection or breakpoints are way harder to
execute. This makes finding the source of an error or explaining unexpected
behaviour harder than on the CPU. Another limitation is the bottleneck
introduced by GPU-CPU readback. In our code the simulation runs on WebGPU, but
rendering uses WebGL via Three.js. It would be more optimal to render directly
using WebGPU, skipping the CPU readback entirely. Our current profiling system
can also be improved to give more detailed information and by using GPU-side
performance counters.

\subsection{Visualization}\label{subsec:visualisation}
% Threejs, instancedmesh, animation?, etc
The visualization is implemented in JavaScript, where our rendering stack is
based on Three.js, which internally uses WebGL. Three.js provides camera
management, geometry and lighting, allowing for higher-level rendering code. We
have used instanced rendering, where each boid is represented by a simple cone
geometry. Moreover, we have added two rectangles to the sides of each cone in order to simulate wings flapping, by making them move continuously up and down in a synchronized manner. The transformation matrices are computed on the GPU and applied per
instance. One of the advantages of using Three.js is the ease of use for rapid
testing and development. Because Three.js abstracts low-level graphics programming, setting up a camera, lighting and basic geometry was easy to do.

We could have improved the visualisation by creating or using Level of Detail
(LOD) systems or frustum culling optimizations. This could have sped up the
rendering part of the simulation even further. By how much is unfortunately not
known, as we do not currently profile the rendering part of the program
explicitely.

\subsection{Benchmarking System}
To accurately quantify the performance of our WebGPU-accelerated simulation, we
developed a benchmarking tool, the \texttt{BoidBenchmarker}. This system is
designed to provide reproducible metrics across different hardware
configurations by controlling the simulation environment and capturing granular
timing data.

\subsubsection{Automated Testing Workflow}
The benchmarking process is divided into distinct operational states to ensure
data integrity:
\begin{itemize}
    \item \textbf{Warm-up Phase:} Upon initiation, the system enters a 10-second \texttt{WARMING\_UP} state. This ensures that the GPU pipelines and caches are stabilized before measurement begins.
    \item \textbf{Recording Phase:} Following the warm-up, the system enters a 10-second \texttt{RECORDING} phase where it captures every frame time, alongside specific simulation and rendering deltas.
    \item \textbf{Environment Control:} To minimize external variance, the benchmarker automatically collapses UI panels, resets the simulation to a clean state, and locks the camera position for the duration of the run.
\end{itemize}

\subsubsection{Metrics and Performance Analysis}
The system collects a variety of metrics to evaluate the efficiency of our
implementation:
\begin{itemize}
    \item \textbf{Timing Breakdown:} The tool records individual durations for the GPU compute passes (simulation) and the Three.js WebGL draw calls (rendering).
    \item \textbf{Stability Metrics:} Beyond average FPS, the system calculates the \say{1\% Low FPS.} This metric is for identifying micro-stuttering that may not be apparent in average frame rate calculations.
    \item \textbf{Hardware Context:} Every report automatically captures the user's OS and hardware information (CPU and GPU) via the \texttt{HardwareInfo} interface to provide context for comparison.
\end{itemize}

\subsubsection{Data Export and Reporting}
Results are processed by the \texttt{PerformanceReportExporter}, which supports
multiple output formats to facilitate analysis and documentation. The system
can generate PNG performance graphs, raw JSON data for historical records, and
pre-formatted LaTeX tables for direct copy-paste support in this report.
Furthermore, the tool supports importing previous JSON results to generate
comparison tables, allowing us to directly measure the impact of different
optimization strategies and different parameter configurations.

\section{Simulating Boid behaviour}\label{sec:boids}
% Basic description of our implementation of boid behaviour
The behaviour of the boids is based on the boids paper by Reynolds~\cite{reynolds1987flocks}. This section describes the algorithm in
general, the modifications we made, and the choice of parameters.

\subsection{The model}\label{subsec:model}

The main model consists of a combination of three steering
behaviours: separation, alignment and cohesion. Separation aims to make sure
boids do not collide or crowd too much. A force is applied opposite to
neighbouring boids, where closer boids contribute more to the force than farther ones. Alignment is
used to steer boids in the same general direction, aligning with their
neighbours. This is calculated by taking the average direction of neighbouring
agents. The cohesion steering behaviour is used to remain a group. The boids
stick together to form a cohesive flock, which can be identified as a group.
Cohesion is calculated by steering towards the average position of neighbouring
boids. Combining these three behaviours with different influence strength
results in a realistic flock of birds. Finding the neighbouring birds is based
on a radius around the bird, as well as a maximum angle of vision directly in front of it. This is used to emulate the
vision of a bird, and any boids outside the range are ignored.

Our implementation consists of a CPU and a GPU part. The CPU part handles the
interaction with the user, which is the definition of parameters, starting and
stopping the simulation, and the mouse cursor interacting with the boids. More
information of mouse interaction in~\ref{subsec:modifications}. After the user
input has been gathered, the current state of the simulation is sent to the
GPU\@. Table~\ref{tab:bindings} displays the data that is sent to the GPU\@.
Note that this is required information for the basic simulation, and does not
include the data needed for any acceleration structures. Section~\ref{sec:accel}
will go into more detail of the chosen structure to find neighbourhoods in a
fast way.

\begin{table}[H]
    \centering
    \begin{tabular}{lll}
        \hline
        Boids      & An array holding boid locations and velocities \\ \hline
        Parameters & The parameters controlling the simulation      \\ \hline
        Matrices   & The output matrices used by the renderer       \\ \hline
    \end{tabular}
    \caption{GPU resources}\label{tab:bindings}
\end{table}

After the simulation has been performed, the altered data is sent back to the
CPU and rendered. The rendering is done as an \texttt{InstancedMesh}\footnote{\texttt{three.js} \href{https://threejs.org/docs/\#InstancedMesh}{documentation on \texttt{InstancedMesh}}}, meaning 
all boids are combined into a single mesh, and drawn as a single draw call.
This is possible, because all boids use the same mesh as a base, and only a
rotation and translation is applied. The usage of an \texttt{InstancedMesh} results in
faster rendering times.

\subsection{Our modifications}\label{subsec:modifications}
While the simulation started of as an exact implementation of Reynold's paper,
we have made some modifications. For instance, we added an interaction
between the boids and the mouse cursor. We have created a system where the boids flee from a
line emitted from the cursor. This system works by casting a ray from the
cursor in the direction of the camera view. The ray origin and direction are
sent to the GPU. On the GPU side, the boid finds the closest point on the line,
and steers away from that point.

Another modification is creating a \say{soft boundary}. This ensure the boids remain
in the scene, but do not create an obvious bounding box shape. This boundary is implemented by adding a small force back to
the center when boids exit the bounding box, which is stronger the further away from the center the model is. They can still exit the cubic bounding area, but are slightly pulled back towards the center. This resulted in
more natural behaviour, while still being able to frame the flocks better.

\subsection{Parameter choice}\label{subsec:parameters}
This section describes the exposed parameters, their function and their default
value. Altering these parameters can change the behaviour, performance or scale
of the simulation. All parameters are listed in Table~\ref{tab:parameters}.

\begin{table}[H]
    \centering
    \begin{tabular}{lp{0.5\linewidth}l}
        \textbf{Parameter}  & \textbf{Function}                                                                                           & \textbf{Default value} \\ \hline
        separation distance & The neighbourhood radius of the separation behaviour                                                        & 25.0                   \\ \hline
        alignment distance  & The neighbourhood radius of the alignment behaviour                                                         & 50.0                   \\ \hline
        cohesion distance   & The neighbourhood radius of the cohesion behaviour                                                          & 50.0                   \\ \hline
        maximum speed       & The maximum units a boid can move in a single timestep                                                      & 5.0                    \\ \hline
        maximum force       & The maximum amount of change to the velocity in a single timestep                                           & 0.1                    \\ \hline
        separation weight   & The relative influence of separation when combining steering behaviours                                     & 1.5                    \\ \hline
        alignment weight    & The relative influence of alignment when combining steering behaviours                                      & 1.0                    \\ \hline
        cohesion weight     & The relative influence of cohesion when combining steering behaviours                                       & 0.5                    \\ \hline
        margin              & The minimum distance to a wall to start steering away from it                                               & 100.0                  \\ \hline
        turnfactor          & The maximum amount of change in velocity to avoid colliding with walls                                      & 0.2                    \\ \hline
        cellsize            & The size of a cell in the spatial hashing datastructure                                                     & 25.0                   \\ \hline
        world max           & The bounding box size of the simulation                                                                     & Calculated             \\ \hline
        grid dim            & The amount of cells in the simulation in the x, y and z direction. The w is used as a total count of cells. & Calculated             \\ \hline
        ray origin          & The origin of the mouse ray. This is the location of the mouse on the near plane of the camera              & Camera pos             \\ \hline
        ray direction       & The direction of the mouse ray. This is the direction of the camera view                                    & Mouse vector           \\ \hline
        flee radius         & The radius of influence of the mouse ray                                                                    & 12\% of width          \\ \hline
    \end{tabular}
    \caption{Parameters}\label{tab:parameters}
\end{table}

\section{Acceleration Structures}\label{sec:accel}
% Describe acceleration structures and ideas we had to optimize our program
% Include ideas we didn't use
% Reference papers or other implementations in .bib file and use \cite to include it in the text

% Please feel free to set a different name to this section, I'm just trying to be somewhat general with the names and we can adapt them later to how the report turns out
% Add more subsections as you please!
To try to get 100,000 boids simulated at a steady 60fps, we have explored a
variety of spatial data acceleration structures. The main bottleneck in the
simulation was finding neighbouring boids. To improve the speed of this, it is
important that the distance calculation is not performed for every boid pair in the simulation.
This would result in a $\mathcal{O}(N^2)$ time complexity. While this is
managable with a low amount of boids, it reduces the fps massively when scaling
to a larger simulation size. This chapter will describe the different data
strucures we have tried, with their strengths and weaknesses.

\subsection{Spatial Hashing}
\label{subsec:spatial_hash}

Spatial hashing is used to accelerate nearest-neighbour queries so each boid
only inspects nearby flockmates instead of the entire set. A regular 3D grid
partitions world space; each boid is mapped to a cell and cells stores linked
lists of boid indices. This reduces the naive $\mathcal{O}(N^2)$ neighbour
search to a local search over the boid population inside a small neighbourhood
(the surrounding $3 \times 3 \times 3$ cell volume).

Implementation details:
\begin{itemize}
    \item \textbf{Grid parameters:} The JavaScript side computes cell size and grid dimensions in \texttt{calculateGridDimensions()}. The chosen \texttt{cell\_size} is based on behaviour distances so neighbours fall into adjacent cells.
    
    \item \textbf{Buffers:} Two GPU storage buffers implement the structure: \texttt{cell\_heads} (an array of atomic integers containing the head index for each cell) and \texttt{boid\_next} (an array of ints forming the linked list next pointers for each boid). Buffers are created and sized in \texttt{initSpatialHashBuffers()}.
    
    \item \textbf{Coordinate helpers:} The compute shader provides \path{pos_to_cell()} and \path{cell_to_index()} helpers to quantize world positions to integer cell coordinates and to flatten 3D coordinates to a 1D cell index.
    
    \item \textbf{Build phase (Pass 1 \& 2):} The GPU executes a two-step build: (a) clear all \texttt{cell\_heads} to \texttt{-1} using \texttt{atomicStore}, and (b) run a parallel insert where each boid computes its cell index and calls \path{atomicExchange(&cell_heads[ci], i32(idx))}. The returned old head is written into \path{boid_next[idx]}, producing a thread-safe linked list per cell.
    
    \item \textbf{Query/update phase (Pass 3):} For each boid the shader computes its home cell, then loops the $3 \times 3 \times 3$ neighbourhood. For each neighbour, cell it reads the head with \path{atomicLoad(\&cell_heads[ci])} and traverses the linked list via \path{boid_next} until \texttt{-1} is reached (the traversal is bounded to avoid pathological loops). During traversal the shader accumulates separation, alignment and cohesion contributions from neighbours.
    
    \item \textbf{All-GPU pipeline:} The implementation keeps the entire update pipeline on the GPU: building the hash, computing forces, and producing per-boid 4x4 transformation matrices for instanced rendering. This avoids expensive CPU $\leftrightarrow$ GPU transfers and enables scaling to many thousands of boids.
\end{itemize}

Key WGSL bindings and shader details are declared near the top of the compute
shader (\path{public/compute-shader.wgsl}): the uniform \path{params} holds
\path{cell_size} and \path{grid_dim}, and storage bindings expose
\path{cell_heads}, \path{boid_next} and the \path{boids} input array. The three
compute passes (clear cells, hash insert, update) are dispatched from
JavaScript in sequence to maintain a deterministic frame update. \bigbreak This
design follows the classic parallel spatial-hash linked-list pattern: atomic
head updates to build lists concurrently, followed by per-boid local
neighbourhood traversals for behaviour computation. It provides a practical
balance between memory use, parallelism and neighbor-search accuracy for
real-time flocking.

\subsection{KD-Tree}
\label{subsec:kdtree}
\todo{kd-tree explanation}

\subsection{Sweep Lines}
\label{subsec:sweepline}
\todo{Sweep lines explanation}

\subsection{Stigmergy}
\label{subsec:stigmergy}

\subsection*{The idea}
The idea behind stigmergy is to use the environment as a medium for indirect
communication between agents. In the context of our boid simulation, this would
involve boids leaving behind some form of \say{pheromone} or marker in the
environment that other boids can sense and react to. This could potentially
allow for more complex and emergent flocking behaviors without requiring direct
communication between boids. \bigbreak

The way Stigmergy is implemented in Tensen's research \cite{order_of_magnitude_boids} is through a texture-based approach, which allows for boids to keep a global interactivity intact throughout the simulation. This contrasts with using Spatial Hashing as described in \ref{subsec:spatial_hash}, which only allows for local interactivity, meaning that boids can only interact with other boids in their local neighborhood. This is because the Spatial Hashing approach relies on a grid-based partitioning of space, which limits the interactions to nearby cells.

\subsubsection*{Our Attempted Implementation}
To explore scaling up our flocking simulation and to test global interactivity, we experimented with integrating concepts from Tensen's work directly into our WebGPU pipeline. Rather than solely relying on standard local interactions, we attempted to map these advanced spatial structures to the GPU.

The core of this experimental optimization was the implementation of a parallel Spatial Hash grid. The intended process was broken down into the following steps on the GPU:
\begin{itemize}
    \item \textbf{Grid Initialization:} The compute shader clears a buffer of cell heads, effectively resetting the spatial grid for the current frame.
    \item \textbf{Spatial Hashing (Linked Lists):} Each boid calculates which 3D grid cell it belongs to based on its current position. Using atomic operations (\texttt{atomicExchange}), the boids concurrently insert themselves into the grid, building a linked list of boids for each cell in a completely parallelized manner.
    \item \textbf{Behavior Update:} During the update pass, boids no longer check every other boid in the simulation. Instead, they determine their current cell and only iterate through the linked lists of the surrounding $3 \times 3 \times 3$ neighboring cells to calculate separation, alignment, and cohesion forces.
    \item \textbf{Matrix Computation:} Finally, the compute shader calculates the $4 \times 4$ transformation matrices for each boid, which are passed directly to Three.js for instanced rendering.
\end{itemize}

\subsubsection*{Motivation for this Exploration}
A naive boid simulation requires every boid to check its distance against every other boid to determine local flockmates. This results in an $\mathcal{O}(N^2)$ time complexity, which severely bottlenecks the simulation on the CPU, typically limiting the flock to just a few thousand individuals before frame rates plummet.
\medbreak
By attempting to implement Tensen's GPU-driven spatial hashing \cite{order_of_magnitude_boids}, we aimed to localize the nearest-neighbour search, vastly reducing the number of distance calculations required per frame. Furthermore, keeping the entire pipeline strictly on the GPU would theoretically eliminate the massive performance overhead of transferring hundreds of thousands of vector states back and forth between CPU and GPU memory.
\medbreak
This architectural choice was intended to bypass traditional computational limits and simulate tens, or even hundreds, of thousands of boids at smooth interactive frame rates, achieving a true order of magnitude performance increase.

\subsubsection*{Challenges and Exclusion from the Final Program}
Ultimately, this approach was excluded from our final program due to persistent implementation issues. Integrating this method caused the boids to become unnaturally attracted to a specific slice of the space where the grid cells were located, forcing them to cluster rigidly in that area rather than moving freely.
\medbreak
We attempted to fix this issue by adjusting the parameters of the spatial hash, but could not find a stable solution. Because of these spatial artifacts, it became impossible to accurately evaluate if this implementation actually resolved the issue of global interactivity. Consequently, we chose to omit this experimental feature from the final deliverable to ensure the stability and expected behavior of the core simulation.

\section{Problems and Limitations}
\label{sec:problems_and_limitations}
% Main problems we faced during the project, limitations of our program/solutions, workarounds we came up with

While our WebGPU-accelerated boid simulation successfully scales to over
100,000 agents, we encountered several notable problems and limitations
regarding debugging, data transfer bottlenecks, and algorithmic artifacts.

\subsection{Debugging Compute Shaders}
One of the primary limitations we ran into was the difficulty of debugging.
Because the WGSL computation happens in parallel on the GPU, standard debugging
methods like variable inspection or setting breakpoints are much harder to
execute. This makes finding the source of an error or explaining unexpected
boid behaviour significantly harder than it would be on the CPU.

\subsection{GPU-CPU Readback Bottleneck}
Another major limitation is the performance bottleneck introduced by GPU-CPU
readback. In our current architecture, the simulation logic runs purely on
WebGPU, but the rendering pipeline uses WebGL via Three.js. This requires
reading the transformation matrices back to the CPU every frame, which limits
maximum performance. It would be much more optimal to render directly using
WebGPU, allowing us to skip the CPU readback entirely.

\subsection{Visualization, Optimization and Profiling}
On the rendering side, we could have improved the visualization by implementing
Level of Detail (LOD) systems or frustum culling optimizations. These additions
could have sped up the rendering part of the simulation even further. However,
implementing these feature would require an overhaul of the rendering pipeline
since all the meshes are instanced meshes we cannot easily modify the LOD
without separating out the mesh instances into groups based on the camera
distance. This would also add additional complexity. Additionally, our overall
profiling system could be improved to provide more detailed information by
utilizing GPU-side performance counters.

\subsection{Algorithmic and Global Interactivity Issues}
Finally, we faced behavioral problems when attempting to implement stigmergy
for global interactivity. Implementing this caused an issue where the boids
became attracted to a specific slice of the space where the grid cells were
located. This made the boids cluster in that area,
preventing them from moving freely through the space. Although we tried to fix
this issue by adjusting the spatial hash parameters, we could not find a
solution that worked well enough. Because of this, it was impossible to
determine if the implementation actually solved the global interactivity
problem, which is why the stigmergy texture approach was excluded from the
final program.

\section{Results}
\label{sec:results}
% Not sure if needed, but somewhere we can show some results of our program, benchmarking, or how far we can push it


%\section{AI prompts}
% Necessary?

\section*{Contributions and Self-reflections}
% Website told us to do it so we each gotta do one

\subsection*{Kian Hamidi}
% Kian's contribution and self-reflection

I entered the project with a strong foundation in WebGPU, Three.js, and flocking algorithms, which allowed me to hit the ground running on the core architecture from the very first week. My primary goal was to ensure the simulation could scale efficiently to meet our ambitious performance targets. To achieve this, I focused heavily on the technical implementation of the spatial hashing logic, which was critical for moving the computational burden of neighbor detection onto the GPU. To validate our progress, I also developed a comprehensive benchmarking and reporting system. This tool allowed us to move beyond anecdotal observations by capturing standardized metrics, such as 1\% low FPS counts and automated LaTeX tables, providing a data-driven backbone for our results section. I also dedicated significant time to researching and prototyping a Stigmergy-based system for global interactivity. Although this remained in the experimental phase due to time constraints and specific visual artifacts—where boids would cluster at the edges of grid slices—the research was invaluable for understanding the limitations and synchronization challenges of our spatial grid.

Beyond the initial implementation, I served as the primary debugger for our WGSL shaders, navigating the complexities of parallel programming to resolve technical hurdles that are often difficult to track in a web environment. Reflecting on the project as a whole, I was incredibly impressed by our team's efficiency. We managed to produce a functional "MVP" very early in the cycle, which significantly reduced the typical end-of-project pressure and gave us the freedom to focus on high-level optimizations and feature expansion. I found it particularly rewarding that every member of the group possessed a high level of general programming fluency. This shared technical competence made merging diverse feature branches and troubleshooting across the codebase a low-friction, collaborative experience. The synergy between our different perspectives, combined with an open line of communication and a collective eagerness to learn, made this one of the most positive team experiences I have had. I am very proud of the cohesive, high-performance system we were able to build together.

\subsection*{Alex Janse}
% Alex's contribution and self-reflection

\subsection*{Jaël Kwakkel}
% Jaël's contribution and self-reflection

In this project, I have worked on the initial setup of the boids simulation,
before optimizing, focussing on the behavior. After that, I researched the
different types of spatial optimization structures including ball tree, octree,
quadtree, spatial hashing and bounding volume hierarchies. On the main code
implementation I have contributed to the spatial hashing implementation,
created the mouse interaction, and created the soft boundary system. Besides
that, I have helped others when we worked together on the programming. On the
report I have mostly written the project description, software section, and the
simulation chapter.

Looking back at this project, I have become better at optimization techniques.
While I know quite a bit about optimization, I have never created spatial
optimization. This was also my first time working with
web-shaders/web-rendering. The different ways to work with shaders, and what
parts to perform on the CPU and GPU was interesting to learn about. In terms of
crowdsimulation I mostly used this project to see how simple rules, or forces,
can lead to complex behaviour and that not a lot is needed to create realistic
movement. Working together with this team was nice. We all had a slightly
different background in programming and simulations and our skills matched up
to create something cool we all contributed to. There are some things I could
have done better. I like to start programming to see what possibilities we
have, but this is not always in the most optimal way. Because I started on a
CPU-only version of the spatial hashing, we had to make some arhitecture
changes to allow for it to run partly on the GPU. In the future I should make
sure how we want to create things before getting stuck with bad achitecture.

\subsection*{Antonio Sato}
% Antonio's contribution and self-reflection
I came to this project with very little experience with JavaScript, three.js and HTML, and no prior experience with GPU programming. For what was relevant for this project, I had only previously worked directly with GLSL in the implementation of vertex and fragment shaders. For this reason, I felt that I wasn't able to do much at the beginning, while the team set up the project, as I still needed to learn the basics. At first, I found myself most useful during our group meetings, when we mostly discussed the overall details of the project, rather than direct implementations, as I was able to find literature on the relevant topics and discuss our goals for the project. 

The group was very efficient during the first couple of weeks, meaning that we already had a simulation running and a functional UI. My main contribution during this part was connecting the modifiable parameters in the UI to the WebGPU pipeline. Further along the line, I attempted implementing an approximate nearest-neighbours search in the GPU \cite{leite2012nearest} in order to optimize the efficiency of the program, but it turned out to not have any meaningful impact on our results. I also experimented with different shapes for the bounding box (e.g. sphere, torus) and added the EXR background to the program, instead of a blank one. Lastly, I implemented a functionality to change vary the color of each boid based on its starting position, as a way to observe how the starting position of boids affected their position and behaviour further into the simulation. I was also responsible for setting up the report and bibliography, and proof-reading.

I believe I could have been more helpful to the team had I communicated better early on in the project and asked for help earlier in how to implement certain systems, as I spent a lot of my time trying to understand the inner-workings of the program until I was able to attempt implementing things. I felt overwhelmed at first, since everyone in the group seemed very knowledgeable of the tools and knew where to start, and once I started properly communicating (about 1.5 to 2 weeks in) most of the basic functionalities had already been implemented. I still feel I was helpful and could contribute during group meetings, and I believe I understand our implementation very well and could reproduce it given enough time, but looking back made me realize I could have started being helpful earlier on.

\break
\nocite{*}
\printbibliography

\end{document}